%% notes-tex/031-unified-representation/sections/7-encoders.tex

\section{Molecular encoders and their coverage\secstatus{todo}}
\label{sec:encoders}

A compound enters the model as a vector, and the choice of encoder is a modeling decision
made on measured coverage and measured behavior rather than on a reading list. Twelve
small-molecule encoders live in \file{torchcell/molecule/}: a \texttt{MoleculeEncoder} turns
a list of SMILES into one fixed-width float32 vector per molecule, raises on a SMILES it
cannot parse rather than returning zeros, and is deterministic across calls. Four are RDKit
fingerprints and descriptors with no learned weights (count and bit ECFP4 at 2{,}048,
feature-class FCFP4, MACCS keys, and the 217 RDKit 2D descriptors); one is Mol2Vec, a
skip-gram over Morgan substructure identifiers; four are SMILES transformers (the two
ChemBERTa-2 heads, MoLFormer-XL, a ZINC RoBERTa); one is Uni-Mol, which embeds an ETKDG
conformer; and one is MolE, the contrastively pretrained graph network built for
antimicrobial activity. Weights that are not on the Hugging Face hub (Mol2Vec, MolE) sit under
\file{\$DATA\_ROOT} with a pydantic manifest recording the source URL, the retrieval command
and the sha256, re-hashed on every load. The registry and its tests are described in the
dendron note \texttt{torchcell.molecule.encoders}.

\notesec{Why these twelve}
Three principles chose the set. The benchmark of 25 pretrained embeddings on 25 datasets
found almost every one statistically indistinguishable from count ECFP \external{arXiv
2508.06199, September 2026 revision}, so fingerprints are the mandatory baseline and every
learned encoder has to earn its place against them. Everything installs from RDKit, pip or
the Hugging Face hub, with one vendored exception. And the four representation families are
each present, so a comparison can say which family carries chemistry-to-response signal rather
than which checkpoint won (Table~\ref{tab:provenance}).
\begin{itemize}
\item \textbf{Fingerprints and descriptors} (count and bit ECFP4, FCFP4, MACCS, the 217
  RDKit descriptors): no training corpus, hashed or curated substructures and
  physicochemical properties \external{Rogers and Hahn 2010, \texttt{10.1021/ci100050t};
  Durant et al.\ 2002, \texttt{10.1021/ci010132r}}. FCFP4 tests whether pharmacophore
  feature classes beat atom identity; the descriptors are what a tolerance biologist
  reasons with (logP, molecular weight, polar surface area).
\item \textbf{Mol2Vec}: skip-gram over Morgan identifiers on 19.9 million ZINC and ChEMBL
  compounds, 300 dimensions, still the metabolomics workhorse \external{Jaeger, Fulle
  and Turk 2018, \texttt{10.1021/acs.jcim.7b00616}}.
\item \textbf{SMILES transformers}: ChemBERTa-2 in its masked-language and
  multitask-regression heads on 77 million PubChem SMILES \external{Ahmad et al.\ 2022,
  arXiv 2209.01712}, MoLFormer-XL on 1.1 billion PubChem and ZINC SMILES with linear
  attention, of which the public checkpoint carries the 10 percent subset \external{Ross
  et al.\ 2022, \texttt{10.1038/s42256-022-00580-7}}, and a community ZINC RoBERTa at
  480 million SMILES added for availability. The two ChemBERTa heads share a backbone,
  which separates SMILES syntax from property supervision.
\item \textbf{3D}: Uni-Mol, pretrained on 209 million conformers by atom masking,
  coordinate denoising and pair-distance prediction \external{Zhou et al.\ 2023, ICLR};
  the one encoder that can say whether conformation matters for the small solvents
  and the phenolics.
\item \textbf{Graph}: MolE, a GIN pretrained with a Barlow Twins objective on 100{,}000
  PubChem structures and validated experimentally for antimicrobial discovery
  \external{Olayo-Alarcon et al.\ 2025, \texttt{10.1038/s41467-025-58804-4}}; whole-cell
  growth inhibition is the published task closest to inhibitor tolerance, and its
  model code is vendored under \file{third\_party/mole/}.
\end{itemize}

%% SOURCE: ENCODER_PROVENANCE in experiments/031-env-chemgen-inhibitor-tolerance/scripts/notes_tex_tables.py
%% (curated metadata, also written to results/encoder_provenance.csv); rendered by
%% notes_tex_representation_tables.py.
\begin{table}[htbp]
\centering
\footnotesize
\caption{Provenance of the twelve encoders: representation family, training corpus,
objective and source. \tcflagext~Temporary: the linked works are not yet in the library
mirror and are given by their public link; the flag lifts when they are mirrored and cited
through the bibliography.}
\label{tab:provenance}
%% GENERATED by experiments/031-env-chemgen-inhibitor-tolerance/scripts/notes_tex_representation_tables.py
%% SOURCE: experiments/031-env-chemgen-inhibitor-tolerance/results/encoder_provenance.csv (ENCODER_PROVENANCE in notes_tex_tables.py, curated metadata; the linked works are not in the library mirror)
%% Do not edit by hand.
\begin{tabular}{p{2.6cm}p{1.9cm}p{4.2cm}p{4.4cm}p{3.2cm}}
\toprule
\textbf{encoder} & \textbf{family} & \textbf{trained on} & \textbf{objective} & \textbf{source} \\
\midrule
ecfp4\_count, ecfp4\_bit & fingerprint & none; hashed radius-2 atom environments & none & \href{https://doi.org/10.1021/ci100050t}{Rogers and Hahn 2010} \\
fcfp4\_count & fingerprint & none; pharmacophore feature invariants & none & \href{https://doi.org/10.1021/ci100050t}{Rogers and Hahn 2010} \\
maccs & fingerprint & none; 166 curated substructure keys & none & \href{https://doi.org/10.1021/ci010132r}{Durant et al. 2002} \\
rdkit\_2d & descriptors & none; 217 physicochemical descriptors & none & \href{https://www.rdkit.org/docs/GettingStartedInPython.html#list-of-available-descriptors}{RDKit 2026.03} \\
mol2vec & fragment embedding & 19.9M ZINC and ChEMBL compounds & skip-gram over Morgan identifiers & \href{https://doi.org/10.1021/acs.jcim.7b00616}{Jaeger, Fulle and Turk 2018} \\
chemberta2\_mlm & SMILES transformer & 77M PubChem SMILES & masked language modeling & \href{https://arxiv.org/abs/2209.01712}{Ahmad et al. 2022} \\
chemberta2\_mtr & SMILES transformer & 77M PubChem SMILES & multitask regression on RDKit properties & \href{https://arxiv.org/abs/2209.01712}{Ahmad et al. 2022} \\
molformer\_xl & SMILES transformer & 1.1B PubChem and ZINC SMILES; public checkpoint is the 10 percent subset & masked language modeling, linear attention & \href{https://doi.org/10.1038/s42256-022-00580-7}{Ross et al. 2022} \\
roberta\_zinc\_480m & SMILES transformer & 480M ZINC SMILES & masked language modeling & \href{https://huggingface.co/entropy/roberta_zinc_480m}{Hugging Face entropy/roberta\_zinc\_480m} \\
unimol\_v1 & 3D & 209M conformers & atom masking, coordinate denoising, pair distances & \href{https://openreview.net/forum?id=6K2RM6wVqKu}{Zhou et al. 2023 (ICLR)} \\
mole\_static & graph (GIN) & 100k PubChem structures & Barlow Twins on masked subgraph views; validated for antimicrobial activity & \href{https://doi.org/10.1038/s41467-025-58804-4}{Olayo-Alarcon et al. 2025} \\
\bottomrule
\end{tabular}

\end{table}

\notesec{Left out}
Chemprop's D-MPNN is trained end to end and belongs to the model phase as an architecture,
not to an embedding comparison. CLAMP, the strongest frozen model in the benchmark, aligns
fingerprints to mammalian and enzymatic assay text and ships only as a research repository;
it is the first candidate for a second pass if the fingerprint arm wins. MolCLR, GraphMVP, 3D
Infomax, Mole-BERT and KPGT are GIN checkpoints of MolE's family in research repositories; GEM
is PaddlePaddle only; CLOOME, MoCoP and InfoAlign encode human Cell Painting or expression
phenotypes whose transfer to yeast growth is untested; MoleculeSTM, MolBERT, R-MAT and CDDD
are covered by the SMILES family already present. Reaction-aware encoders were left out on
purpose: a reaction-hypergraph embedding places a metabolite by the reactions it takes part
in, which is the right context for a Yeast9 metabolite and no context at all for an inhibitor
whose fate in the cell is not in the model. The metabolites among the medium components can
carry a reaction branch (Section~\ref{sec:metabolites}); the inhibitors cannot.

\notesec{Two caveats for these molecules}
The corpora are drug-like, 20 to 40 heavy atoms. Ethanol and isobutanol have 3 and 5, so
ECFP is nearly empty and a transformer embedding is poorly determined for them, which is part
of why isobutanol has no close chemical neighbor in any partner (Section~\ref{sec:chemistry}).
The salts (sodium acetate, the imidazolium chlorides, the metal chlorides) are multi-fragment
SMILES that fingerprints handle and Uni-Mol cannot conformer.

\notesec{Coverage over the five datasets}
Table~\ref{tab:coverage} gives every encoder over the dosed compounds of all five datasets,
5{,}472 distinct InChIKeys in the union. Every one of them now has a structure, and that
sentence hides a fix worth recording. The curated identity table has a row for every
Vanacloig, Hillenmeyer and Wildenhain compound, but Hoepfner derives 88 of its 148 InChIKeys
from the SMILES released in its Table S1 and those keys were never curated, so until the
served record's own SMILES was read as the structure source those 88 compounds, 60 percent of
Hoepfner's chemistry, were invisible to every encoder. The precedence is the curated table
where it has a row and the served record where it has none, because the two routes measurably
disagree on the stereo block for some compounds. Eleven of the twelve encoders embed the whole
union. Uni-Mol refuses the compounds for which no three-dimensional conformer exists, the
isolated ions and metal salts, each named in \file{results/embeddings/failures.json}; its
wrapper library would otherwise substitute flat or all-zero coordinates and serve a vector
from nothing, which is why the encoder's \texttt{check} reruns the conformer step and raises.
One dosed Hillenmeyer name, tunicamycin, is a homologue mixture with no InChIKey and is not a
compound here.

%% SOURCE: experiments/031-env-chemgen-inhibitor-tolerance/scripts/embed_compounds.py ->
%% results/embedding_coverage.csv, rendered by notes_tex_representation_tables.py.
\begin{table}[htbp]
\centering
\footnotesize
\caption{Coverage of every encoder over the dosed compounds of the five datasets (distinct
InChIKeys with a dosed record). The last column is wall time in seconds to embed the 5{,}472
compounds of the union on one shared GPU, model load excluded.}
\label{tab:coverage}
%% GENERATED by experiments/031-env-chemgen-inhibitor-tolerance/scripts/notes_tex_representation_tables.py
%% SOURCE: experiments/031-env-chemgen-inhibitor-tolerance/results/embedding_coverage.csv from embed_compounds.py
%% Do not edit by hand.
\begin{tabular}{lrrrrrrr}
\toprule
\textbf{encoder} & \textbf{dim} & \textbf{Vanacloig} & \textbf{HOM} & \textbf{HET} & \textbf{Hoepfner} & \textbf{Wildenhain} & \textbf{s} \\
\midrule
chemberta2\_mlm & 384 & 41 of 41 & 114 of 114 & 290 of 290 & 148 of 148 & 5170 of 5170 & 3.4 \\
chemberta2\_mtr & 384 & 41 of 41 & 114 of 114 & 290 of 290 & 148 of 148 & 5170 of 5170 & 3.4 \\
ecfp4\_bit & 2,048 & 41 of 41 & 114 of 114 & 290 of 290 & 148 of 148 & 5170 of 5170 & 1.3 \\
ecfp4\_count & 2,048 & 41 of 41 & 114 of 114 & 290 of 290 & 148 of 148 & 5170 of 5170 & 1.4 \\
fcfp4\_count & 2,048 & 41 of 41 & 114 of 114 & 290 of 290 & 148 of 148 & 5170 of 5170 & 1.8 \\
maccs & 167 & 41 of 41 & 114 of 114 & 290 of 290 & 148 of 148 & 5170 of 5170 & 4.9 \\
mol2vec & 300 & 41 of 41 & 114 of 114 & 290 of 290 & 148 of 148 & 5170 of 5170 & 1.9 \\
mole\_static & 1,000 & 41 of 41 & 114 of 114 & 290 of 290 & 148 of 148 & 5170 of 5170 & 4.9 \\
molformer\_xl & 768 & 41 of 41 & 114 of 114 & 290 of 290 & 148 of 148 & 5170 of 5170 & 8.0 \\
rdkit\_2d & 217 & 41 of 41 & 114 of 114 & 290 of 290 & 148 of 148 & 5170 of 5170 & 41.5 \\
roberta\_zinc\_480m & 768 & 41 of 41 & 114 of 114 & 290 of 290 & 148 of 148 & 5170 of 5170 & 9.5 \\
unimol\_v1 & 512 & 41 of 41 & 103 of 114 & 283 of 290 & 136 of 148 & 5163 of 5170 & 807.3 \\
\bottomrule
\end{tabular}

\end{table}

\notesec{Cost}
The fingerprints, Mol2Vec, MolE and the SMILES transformers embed the union in seconds; the
RDKit descriptors take longer only because of a few slow descriptors; Uni-Mol takes minutes
because of conformer generation. At the scale of these datasets no encoder is expensive, and
the choice is made on what the vectors do. The same twelve encoders also cover the 37 medium
components and the 894 structure-bearing Yeast9 species (Section~\ref{sec:metabolites}), so
one embedding table serves every molecular slot in the representation.

\notesec{Chemical similarity against response similarity}
Table~\ref{tab:chemsim} measures, for every encoder and both Hillenmeyer arms, whether two
conditions that are chemically similar respond similarly across datasets: the Spearman
between chemical similarity (Tanimoto for fingerprints, cosine for dense embeddings) and the
cross-dataset response Spearman over every Vanacloig-by-partner condition pair. Against the
homozygous arm the correlation is 0.08 to 0.09 for count and bit ECFP4, MoLFormer-XL and
Mol2Vec, 0.05 to 0.06 for the ChemBERTa-2 heads, and absent for MolE, the RDKit descriptors
and the ZINC RoBERTa; against the heterozygous arm every value is below 0.04. The effect is
real and small, because almost every pair sits at a Tanimoto below 0.25 where response
similarity is centered on zero. The cosine medians of the dense encoders, 0.8 to 0.95, are not
comparable to Tanimoto: cosine between mean-pooled transformer states is high for any two
molecules, so only the fingerprint rows say how close the chemistry is.

%% SOURCE: experiments/031-env-chemgen-inhibitor-tolerance/scripts/chemical_similarity.py ->
%% results/chemical_similarity_summary.csv, rendered by notes_tex_representation_tables.py.
\begin{table}[htbp]
\centering
\footnotesize
\caption{Per encoder and Hillenmeyer arm: the median similarity of each Vanacloig compound
to its nearest partner compound, the Spearman between chemical similarity and cross-dataset
response Spearman over every condition pair with its $p$ value, and the median response
Spearman of the top decile of pairs by chemistry against the rest.}
\label{tab:chemsim}
%% GENERATED by experiments/031-env-chemgen-inhibitor-tolerance/scripts/notes_tex_representation_tables.py
%% SOURCE: experiments/031-env-chemgen-inhibitor-tolerance/results/chemical_similarity_summary.csv from chemical_similarity.py
%% Do not edit by hand.
\begin{tabular}{llrrrrr}
\toprule
\textbf{encoder} & \textbf{partner} & \textbf{NN median} & \textbf{$\rho$(chem, response)} & \textbf{$p$} & \textbf{top-decile $\rho$} & \textbf{rest $\rho$} \\
\midrule
chemberta2\_mlm & HET & 0.85 & 0.037 & 4.8e-05 & 0.005 & 0.002 \\
chemberta2\_mtr & HET & 0.94 & 0.031 & 6.7e-04 & 0.005 & 0.002 \\
ecfp4\_bit & HET & 0.33 & 0.014 & 1.4e-01 & 0.003 & 0.002 \\
ecfp4\_count & HET & 0.32 & 0.033 & 3.3e-04 & 0.004 & 0.002 \\
fcfp4\_count & HET & 0.44 & 0.039 & 2.6e-05 & 0.005 & 0.002 \\
maccs & HET & 0.60 & 0.002 & 8.6e-01 & 0.003 & 0.002 \\
mol2vec & HET & 0.94 & 0.015 & 9.8e-02 & 0.003 & 0.002 \\
mole\_static & HET & 0.93 & -0.005 & 5.6e-01 & 0.002 & 0.002 \\
molformer\_xl & HET & 0.82 & 0.040 & 1.1e-05 & 0.006 & 0.002 \\
rdkit\_2d & HET & 0.72 & 0.036 & 7.2e-05 & 0.005 & 0.002 \\
roberta\_zinc\_480m & HET & 0.85 & 0.011 & 2.1e-01 & 0.003 & 0.002 \\
unimol\_v1 & HET & 0.90 & 0.033 & 3.4e-04 & 0.004 & 0.002 \\
chemberta2\_mlm & HOM & 0.85 & 0.064 & 1.3e-05 & 0.019 & 0.011 \\
chemberta2\_mtr & HOM & 0.93 & 0.054 & 2.4e-04 & 0.015 & 0.011 \\
ecfp4\_bit & HOM & 0.33 & 0.089 & 1.3e-09 & 0.011 & 0.012 \\
ecfp4\_count & HOM & 0.33 & 0.080 & 4.7e-08 & 0.014 & 0.011 \\
fcfp4\_count & HOM & 0.39 & 0.038 & 9.8e-03 & 0.010 & 0.012 \\
maccs & HOM & 0.55 & 0.041 & 5.4e-03 & 0.014 & 0.011 \\
mol2vec & HOM & 0.95 & 0.078 & 8.8e-08 & 0.020 & 0.011 \\
mole\_static & HOM & 0.91 & 0.008 & 5.7e-01 & 0.011 & 0.012 \\
molformer\_xl & HOM & 0.81 & 0.079 & 7.0e-08 & 0.022 & 0.011 \\
rdkit\_2d & HOM & 0.71 & 0.016 & 2.7e-01 & 0.012 & 0.011 \\
roberta\_zinc\_480m & HOM & 0.83 & -0.028 & 5.6e-02 & 0.009 & 0.012 \\
unimol\_v1 & HOM & 0.90 & 0.042 & 6.4e-03 & 0.012 & 0.011 \\
\bottomrule
\end{tabular}

\end{table}

Read together with Section~\ref{sec:established}, the encoders divide by where they are
asked to work. Within Vanacloig, where medium, dose basis and readout are constant, a ridge
map from a fingerprint reaches 0.31 on the compound-specific target. Across datasets the same
fingerprints reach 0.09. No encoder in the set closes that gap, which is the measured reason
the representation carries the dataset and dose-basis tokens rather than asking chemistry to
bridge regimes on its own.
